\begin{pastexam}[label={px:gw3-098-ds-04}]{關務三等資訊處理 98 年 第四題}
（申論；20 分）

今有一採用開放位址（open addressing）方式儲存鍵值，大小為 14 的雜湊表（hash table），其雜湊函數為
$H_i(k) = h_1(k) + i\,h_2(k) \pmod{14}$，$i = 0, 1, \ldots, 13$，其中 $h_1(k) = k \pmod{14}$。依序存入下列鍵值
15、18、42、19、10、28、27、38、80、55。
\begin{enumerate}[leftmargin=2.6em, labelsep=0.4em]
\item[㈠] 設 $h_2(k) \equiv 1$，請列出最後雜湊表的結果與使用雜湊函數的次數。（10 分）
\item[㈡] 設 $h_2(k) = 1 + k \pmod{13}$，請列出最後雜湊表的結果與使用雜湊函數的次數。（10 分）
\end{enumerate}
\end{pastexam}

\begin{pastsoln}
「使用雜湊函數的次數」解讀為計算 $H_i(k)$ 的次數，也就是探測次數：每個鍵從 $i = 0$ 算起，位址被占用就把 $i$ 加 $1$
再算一次，直到找到空位（資結書第~12.3~節）。 ㈠的 $h_2 \equiv 1$ 就是線性探測，㈡是雙重雜湊。

各鍵的 $h_1(k) = k \bmod 14$：$15 \to 1$、$18 \to 4$、$42 \to 0$、$19 \to 5$、$10 \to 10$、$28 \to 0$、$27 \to 13$、
$38 \to 10$、$80 \to 10$、$55 \to 13$。

\textbf{㈠ $h_2(k) = 1$（線性探測）}：
\begin{center}
\small
\begin{tabular}{clc}
\toprule
鍵 & 探測的位址（$i = 0, 1, \ldots$） & 次數 \\
\midrule
15 & 1 & 1 \\
18 & 4 & 1 \\
42 & 0 & 1 \\
19 & 5 & 1 \\
10 & 10 & 1 \\
28 & 0（42）、1（15）、\textbf{2} & 3 \\
27 & 13 & 1 \\
38 & 10（10）、\textbf{11} & 2 \\
80 & 10（10）、11（38）、\textbf{12} & 3 \\
55 & 13（27）、0（42）、1（15）、2（28）、\textbf{3} & 5 \\
\bottomrule
\end{tabular}
\end{center}
（括號內是占用該位址的鍵。 位址 $13$ 的下一個是 $(13 + 1) \bmod 14 = 0$，繞回開頭。）
\begin{center}
\small
\setlength{\tabcolsep}{4pt}
\begin{tabular}{c|cccccccccccccc}
位址 & 0 & 1 & 2 & 3 & 4 & 5 & 6 & 7 & 8 & 9 & 10 & 11 & 12 & 13 \\
\hline
鍵 & 42 & 15 & 28 & 55 & 18 & 19 & & & & & 10 & 38 & 80 & 27 \\
\end{tabular}
\end{center}
使用雜湊函數共 $1 + 1 + 1 + 1 + 1 + 3 + 1 + 2 + 3 + 5 = 19$ 次。 可以看到線性探測的「群聚」：位址 $10$--$13$、$0$--$5$
連成一大段，55 要走過一整段才找到空位。

\textbf{㈡ $h_2(k) = 1 + (k \bmod 13)$（雙重雜湊）}：前五個鍵都沒有碰撞，各 $1$ 次。 之後：
\begin{itemize}
\item $28$：$h_1 = 0$（42 占用）；$h_2 = 1 + (28 \bmod 13) = 1 + 2 = 3$；$H_1 = 0 + 3 = 3$，空，放入。 $2$ 次。
\item $27$：$h_1 = 13$，空，放入。 $1$ 次。
\item $38$：$h_1 = 10$（10 占用）；$h_2 = 1 + (38 \bmod 13) = 1 + 12 = 13$；$H_1 = (10 + 13) \bmod 14 = 9$，空，放入。 $2$ 次。
\item $80$：$h_1 = 10$（占用）；$h_2 = 1 + (80 \bmod 13) = 1 + 2 = 3$；$H_1 = 13$（27 占用）；$H_2 = (10 + 6) \bmod 14 = 2$，空，放入。 $3$ 次。
\item $55$：$h_1 = 13$（27 占用）；$h_2 = 1 + (55 \bmod 13) = 1 + 3 = 4$；$H_1 = 17 \bmod 14 = 3$（28 占用）；
$H_2 = 21 \bmod 14 = 7$，空，放入。 $3$ 次。
\end{itemize}
\begin{center}
\small
\setlength{\tabcolsep}{4pt}
\begin{tabular}{c|cccccccccccccc}
位址 & 0 & 1 & 2 & 3 & 4 & 5 & 6 & 7 & 8 & 9 & 10 & 11 & 12 & 13 \\
\hline
鍵 & 42 & 15 & 80 & 28 & 18 & 19 & & 55 & & 38 & 10 & & & 27 \\
\end{tabular}
\end{center}
使用雜湊函數共 $1 + 1 + 1 + 1 + 1 + 2 + 1 + 2 + 3 + 3 = 16$ 次，比線性探測少，因為碰撞的鍵各自用不同的步長跳開，
不會擠在同一段。

補充：表的大小 $14$ 不是質數，$h_2(k)$ 若和 $14$ 有公因數（例如 $h_2(18) = 1 + 5 = 6$），探測序列只會走到一部分的
位址（$\gcd(6, 14) = 2$，只走 $7$ 個），表還沒滿也可能找不到空位（資結書定理~12.4 要求 $\gcd(h_2(k), m) = 1$）。
這一題的鍵碰巧都在前三次內找到空位，沒有出問題；實務上雙重雜湊的表大小應取質數。
\end{pastsoln}
